\title{Byte Latent Transformer: Patches Scale Better Than Tokens}

\begin{abstract}
We present the Byte Latent Transformer ({\textsc BLT}), a novel byte-level LLM framework that, for the first time, achieves parity with tokenization-based LLMs at scale, while also delivering notable gains in inference speed and model robustness.
{\textsc BLT} operates by encoding byte sequences into patches of varying lengths, which become the core units of computation within the model. These patches are delineated dynamically, using the entropy of upcoming bytes to allocate additional computational and model resources precisely where the input data is more intricate. We conduct the initial scaling analysis of byte-level language models constrained by floating point operations, exploring up to 8B parameters and 4T bytes during training. Our findings verify that large-scale training directly on raw bytes—without dependence on predetermined vocabularies—is practical and effective. By favoring longer patches in predictable regions, both training and inference become more efficient, and the model demonstrates superior qualitative abilities in reasoning and handling atypical scenarios. Ultimately, {\textsc BLT} achieves markedly improved scaling over token-based models at equivalent inference costs by simultaneously expanding patch length and model complexity.
\end{abstract}

\section{Introduction}
\label{section:intro}

We present the Byte Latent Transformer~(\textbf{BLT}), a novel architecture that dispenses with tokenization by directly processing raw byte sequences. For the first time, our approach achieves performance on par with models utilizing conventional tokenization techniques at large scales, while delivering notable gains in efficiency and robustness (\S\ref{sec:robustness}). In contrast, contemporary large language models (LLMs) are trained in an almost end-to-end fashion except for a heuristic tokenization stage—this step converts bytes into a predetermined vocabulary of tokens. The choice and structure of these tokens introduce biases during sequence compression and contribute to well-known limitations, including heightened sensitivity to domains and modalities~\citep{dagangetting}, increased vulnerability to input perturbations~(\S\ref{sec:robustness}), insufficient modeling of orthographic patterns~\citep{edman2024cute}, and disparities in multilingual handling~\citep{liang2023xlm,petrov2024language,limisiewicz2024myte}.

Previously, tokenization has played a vital role, as training llms directly on bytes incurs excessive computational expense at scale because of the resultant lengthy input sequences~\citep{xue2022byt5}. Earlier studies address this challenge by introducing more efficient self-attention mechanisms~\citep{el2020characterbert, clark2022canine} or adopting architectures that forgo attention entirely~\citep{wang2024mambabyte}~(\S\ref{section:related_work}). Nonetheless, these approaches are primarily beneficial for the training of \textit{small models}. When scaling up, most of the computational burden in a Transformer arises from the substantial feed-forward layers that must process each byte, rather than from the attention computations themselves.

We introduce a dynamic and learnable approach for organizing bytes into \textit{patches}~(\S\ref{section:patching}), complemented by a novel model architecture that integrates both byte-level and patch-level information. In contrast to tokenization, {\textsc BLT} does not rely on a predefined vocabulary for patches. Instead, any selection of bytes can be encoded into latent patch representations through compact, learnable encoder and decoder components. Our findings demonstrate that this strategy enables a \textit{more} effective distribution of computational resources compared to models utilizing tokenization.

Tokenization-based LLMs dedicate an equal amount of computational resources to each token. This design prioritizes simplicity but can compromise efficiency, since token formation relies on compression strategies that do not consistently reflect the actual difficulty of prediction. The foundation of our architecture rests on the principle that computational effort should be adaptively assigned according to necessity. For instance, complex models like large transformers are typically redundant when predicting common word endings, which are relatively straightforward and entail low uncertainty, unlike the more challenging task of selecting the first word in a sentence. This principle is embodied in the structure of {\textsc BLT}~(\S\ref{section:architecture}), which incorporates three distinct transformer modules: two compact, byte-level \textit{local models}, and a substantially larger global \textit{latent transformer}~(\autoref{fig:arch_high_level}). 

To facilitate dynamic computation assignment and the grouping of bytes into patches, {\textsc BLT} organizes data according to the entropy associated with next-byte predictions. This method produces context-aware clusters of bytes that maintain relatively consistent information density.

We introduce the first comprehensive flop-controlled scaling analysis of byte-level models scaling up to 8B parameters and exposed to 4T training bytes, demonstrating the feasibility of training end-to-end directly from raw bytes without relying on predetermined vocabulary tokenization. Notably, {\textsc BLT} attains training flop-controlled results that are comparable to Llama 3, yet achieves up to 50\% reduction in inference flops~(\S\ref{section:scaling}).\footnote{Here, computational cost is measured as the count of Floating Point OPerations (flops) performed.} Furthermore, operating natively on byte sequences yields clear advancements in modeling rare and less frequent patterns in data. {\textsc BLT} models exhibit greater resilience to noisy text compared to those relying on tokenizers, as well as superior capacity for character-level interpretation across tasks involving orthographic awareness, phonological properties, and machine translation in low-resource scenarios~(\S\ref{sec:robustness}). 

Moreover, {\textsc BLT} architectures support simultaneous scaling of both model and patch sizes while adhering to fixed inference flop constraints. Increasing patch size generally results in computational savings, which can instead be allocated to enlarging the global latent transformer, since it is invoked less frequently overall. Through a series of inference-flop controlled scaling studies~(\autoref{fig:fixed_inference_scaling}), we report substantially more favorable scaling behaviors relative to conventional tokenization-based models.

To conclude, this work offers the following key contributions: 1) We present {\textsc BLT}, a byte latent LLM framework that adaptively allocates computational resources to enhance FLOP efficiency; 2) We provide evidence that our approach achieves parity with Llama 3 in terms of training FLOP control up to the 8B parameter scale, with the ability to exchange slight reductions in evaluation scores for up to a 50\% improvement in FLOP efficiency; 3) {\textsc BLT} enables a novel approach to LLM scaling, allowing model capacity to increase without exceeding a predetermined inference computation limit; 4) We illustrate that {\textsc BLT} architectures display greater resilience to input perturbations and exhibit finer sensitivity to sub-word structures in data—features often overlooked by conventional token-based LLMs. The source code for both training and inference with {\textsc BLT} is available at~\url{https://github.com/facebookresearch/blt}.

\section{Patching: From Individual Bytes to Groups of Bytes}
\label{section:patching}

Dividing a byte sequence into \textit{patches} enables {\textsc BLT} to flexibly adjust computational effort according to the surrounding context. Figure~\ref{fig:patching_types} presents multiple strategies for partitioning bytes into patches. More precisely, a patching function $f_p$ divides a byte sequence $\pmb{x}=\{x_i,|i=1,\ldots n\}$ of length $n$ into $m$ patches $\pmb{p}=\{p_j|j=1,\ldots,m\}$, with $m < n$, by assigning each $x_i$ a label in the set \{0,1\}, where a value of 1 signals the initiation of a new patch. For both patch-based and token-based approaches, the principal driver of computational complexity is the number of main Transformer operations, which in {\textsc BLT} directly corresponds to the count of patches required to encode the input as determined by the patching function. Therefore, the \textit{patch size}—the expected length of each patch—emerges as the chief determinant of resource usage for training and inference for a chosen patching strategy~(\S\ref{section:flops}).

In what follows, we introduce three patching strategies: grouping a fixed count of bytes per patch~(\S\ref{section:static-patch}), segmenting at whitespaces~(\S\ref{section:white-patch}), and using entropy estimates from a compact byte language model to generate variable-length patches dynamically~(\S\ref{section:dyn-patch}). We conclude with a discussion on incremental patching and the distinctions between patching and tokenization~(\S\ref{section:bpe-patch}).

\subsection{Strided Patching Every K Bytes}
\label{section:static-patch}

One of the simplest methods for grouping bytes is to divide them into patches of a fixed length $k$, as employed in MegaByte~\citep{yu2023megabyte}. This approach, which uses a constant stride, is straightforward to implement during both training and inference, allows for easy adjustment of the average patch size, and simplifies the management of computational costs (in flops). Despite its practical advantages, there are considerable drawbacks. Chiefly, computational resources are not adaptively focused where they are most beneficial: transformer steps may be squandered on trivial content like whitespace in code, while regions dense with meaningful information—such as mathematical notation—may not receive adequate attention. Additionally, this static patching method results in uneven and non-adaptive segmentation of byte sequences, causing issues like identical words being divided inconsistently across contexts.

\subsection{Space Patching}
\label{section:white-patch}

\citet{slagle2024spacebyte} introduces an enhancement to strided patching by generating new patches following any space-like bytes\footnote{
    Space-like bytes are defined as any byte that is not a latin character, digit, or utf-8 continuation byte. Additionally, each patch is required to include at least one non space-like byte. }—positions that commonly serve as boundaries for words or other linguistic elements in various languages. Under the Space patching approach, each word is allocated a latent transformer step (increasing computational cost) to enable precise modeling of every word boundary. This design ensures that words receive consistent patching treatment across sequences and dedicates computational resources to challenging predictions, particularly those that occur just after spaces. For instance, inferring the initial byte of the response to ``Who composed the Magic Flute?\uline{\hspace{.6em}}'' is considerably more challenging than predicting the subsequent bytes after “M”, as the initial character dramatically restricts the set of plausible answers, thereby making the remainder of the word (e.g., “Mozart”) easier to predict. Despite its strengths, space patching is limited in its applicability to certain languages or textual domains and, critically, lacks flexibility to adapt patch sizes. In the following section, we propose an alternative patching technique that leverages the observation that the initial bytes of words tend to present the greatest prediction difficulty, while inherently providing the ability to control patch size.

\subsection{Entropy Patching: Using Next-Byte Entropies from a Small Byte LM}
\label{section:dyn-patch}

Instead of utilizing rule-based heuristics like whitespace, we propose a data-driven method to detect points of elevated uncertainty in next-byte predictions. Specifically, we present \textit{entropy patching}, a technique that leverages entropy measurements to determine patch boundaries.

We train a small byte-level auto-regressive language model on the training data for {\textsc BLT} and compute next byte entropies under the LM distribution $p_{e}$ over the byte vocabulary $\mathcal{V}$:

\begin{equation}
H(x_i) = \sum_{v \in \mathcal{V}} p_{e}(x_i=v|\pmb{x}_{<i}) \log p_{e}(x_i=v|\pmb{x}_{<i})
\end{equation}

We explore two approaches for detecting patch boundaries based on the entropies $H(x_i)$. The initial method selects points exceeding a predetermined global entropy threshold, as depicted in~\autoref{fig:patching}. The alternative approach looks for points where the entropy is relatively elevated compared to the preceding value. This latter strategy effectively highlights points that interrupt an otherwise approximately monotonically decreasing entropy trend within a patch.

\begin{align*}
    \text{Global Constraint}\hspace{1cm}H(x_t)& > \theta_g\\
    \text{Approx. Monotonic Constraint}\hspace{1cm}H(x_t)& - H(x_{t-1}) > \theta_r
\end{align*}

Patch boundaries are determined through a simple preprocessing procedure that takes place as part of the dataloading process. This approach contrasts with the method proposed by~\citet{nawrot-etal-2023-efficient}, who employ a trained classifier to detect entropy-based patch boundaries. Within our experiments~(\S\ref{section:experiments}), we evaluate and contrast these two strategies for differentiating between bytes of low and high entropy.

\subsection{The Byte-Pair Encoding (BPE) Tokenizer and Incremental Patching}
\label{section:incremental-patching}
\label{section:bpe-patch}

A large number of contemporary llms, such as our baseline Llama 3, utilize subword tokenizers, for instance bpe~\citep{gage1994new, sennrich-etal-2016-neural}. Throughout this paper, we designate ``tokens'' as byte groupings taken from a \textit{finite} vocabulary that is set before training begins, in contrast to ``patches,'' which denote flexible, dynamically created groupings that lack a predetermined vocabulary. An essential distinction separating tokens from patches is that, in the case of tokens, the model does not possess explicit access to the original byte-level information.

An important advancement offered by {\textsc BLT} compared to models that rely on tokenization lies in its reconfiguration of the balance between vocabulary size and computational demands. In conventional llms, enlarging the vocabulary typically results in tokens that encompass more bytes on average, thereby reducing the number of decoding steps required; however, this comes at the expense of a larger output dimensionality for the model’s final projection layer. This fundamental trade off constrains tokenization-based methods from achieving notable flexibility in both token granularity and inference efficiency. To illustrate, Llama 3 manages to boost its average token length from 3.7 to 4.4 bytes, but this improvement necessitates a fourfold expansion of its embedding table compared to that of Llama 2.

During generation, {\textsc BLT} must determine at each step of the byte sequence whether a patch boundary has been reached, as this affects whether additional computation is triggered using the Latent Transformer. Importantly, this determination must be made without relying on any bytes that have yet to be generated; that is, patch segmentation decisions cannot use information from future portions of the byte sequence. More formally, we define a patching scheme $f_p$ as incremental if it meets the following criterion:
$$f_p(\pmb{x}_{<i}) = f_p(\pmb{x})_{<i}$$
Notably, bpe does not satisfy the requirements of incremental patching, since the same initial sequence can be split into tokens differently depending on subsequent bytes, thus violating the condition stated above.\footnote{Inserting a special delimiter token to mark patch boundaries could make bpe incremental; however, this would increase the total number of bytes in the sequence.}

\section{{\textsc BLT} Architecture}
\label{section:architecture}
The {\textsc BLT} framework consists of three components: a global autoregressive language model, which processes sequences of patch-level representations, and two lightweight local models responsible for mapping byte sequences to patch representations and for reconstructing the original bytes from patch representations, respectively~(\autoref{fig:arch_high_level}).

\subsection{Latent Global Transformer Model}

\textit{The Latent Global Transformer} refers to an autoregressive transformer architecture $\mathcal{G}$ composed of $l_{\mathcal{G}}$ layers, designed to transform sequences of latent input patch embeddings $p_j$ into corresponding sequences of output patch embeddings $o_j$. Within this work, indices $j$ and $i$ are reserved for denoting patches and bytes, respectively. For the global component, a block-causal attention mask~\citep{dubey2024llama} is employed, ensuring that each patch only attends to itself and all preceding patches within a single document. Because this transformer accounts for most of the computational workload both during pre-training and inference, strategically deciding when to apply it provides a mechanism for adapting the computation required for particular segments of input and output depending on their respective complexities.

\subsection{Local Encoder}
\label{section:local}

The \textit{Local Encoder Model}, represented by $\mathcal{E}$, utilizes a compact transformer architecture characterized by a significantly smaller layer count ($l_{\mathcal{E}} << l_{\mathcal{G}}$) in comparison to the global model. Its principal function is to effectively encode an input byte sequence $b_i$ into high-dimensional patch representations $p_j$. Notably, a distinguishing modification from standard transformer configurations is the insertion of a cross-attention mechanism subsequent to each transformer layer; this mechanism aggregates the byte-level features into the corresponding patch representations, as illustrated in~(\autoref{fig:crossattn}).

The processing pipeline commences by embedding the input bytes $b_i$ through a lookup in a matrix of dimension $\mathbb{R}^{256 \times h_{\mathcal{E}}}$, producing vectors $x_i$. This step may incorporate supplementary information via hash-embeddings, elaborated in~(\S\ref{sec:hash-emb}). The embedded inputs then proceed through multiple alternating layers consisting of transformer operations interleaved with cross-attention blocks~(\S\ref{sec:enc-cross-attn}), incrementally evolving the representations into patch-level embeddings $p_i$ that are subsequently forwarded to the global transformer model, $\mathcal{G}$.

Within the transformer blocks, attention is constrained via a \textit{local block causal} mask, such that each token is permitted to attend to a preceding fixed-size context window of $w_{\mathcal{E}}$ bytes. While this window can extend over patch boundaries if necessary, it is strictly contained within the boundaries of the document. The subsections that follow elaborate on the technical particulars of both the embedding methods and the cross-attention modules.

\subsubsection{Encoder Hash n-gram Embeddings}
\label{sec:ngram-emb}
\label{sec:hash-emb}
An essential factor in generating strong and meaningful representations at every step $i$ is embedding contextual information from earlier bytes. In {\textsc BLT}, this is accomplished by jointly encoding the byte $b_i$ on its own and within the context of a byte n-gram. For each position $i$, the method begins by forming byte-grams

\begin{align}
    g_{i,n}=\{b_{i-n + 1},\ldots, b_i\}
\end{align}

for every byte position $i$ and for each $n$ in the range from three to eight.\footnote{
Byte-grams whose length $n$ exceeds the position $i$ ($i < n$) are excluded from consideration. }

Next, we propose hash $n$-gram embeddings, wherein every possible byte $n$-gram is transformed using a hash function to produce an index into an embedding table $E_{n}^{hash}$ of a predetermined size, defined separately for each $n \in\{3,4,5,6,7,8\}$~\citep{bai2010learning}. 
The corresponding embedding is summed with that of the byte itself, and this total is subsequently normalized before being submitted to the local encoder component. The computation of the enhanced embedding proceeds as follows:

\begin{align}
e_i &= x_i + \sum_{n = 3,...,8}E_{n}^{hash}(\text{Hash}(g_{i,n})) \\
\text{where, Hash}(g_{i,n}) &= \text{RollPolyHash}(g_{i,n}) \% |E_{n}^{hash}|
\end{align}

We scale $e_i$ by dividing by the total number of $n$-gram sizes plus one, and employ RollPolyHash as specified in Appendix~\ref{appendix:rollpolyhash}. Section~\ref{section:ablations} presents an ablation study exploring the influence of $n$-gram hash embeddings across various choices of $n$ and embedding table capacities, with respect to trends in flop-controlled scaling laws. Beyond utilizing hash-based $n$-gram embeddings, we conducted additional investigations using frequency-based $n$-gram embeddings; further information about these experiments is available in Appendix~\ref{appendix:ngrams}.

\subsubsection{Encoder Multi-Headed Cross-Attention}
\label{sec:enc-cross-attn}
Our approach largely mirrors the cross-attention mechanism applied to the input in the Perceiver architecture~\citep{jaegle2021perceiver}. A primary distinction in our framework is that latent representations are tied to the variable patch representations, rather than arising from a static collection of latent variables~(\autoref{fig:crossattn}). Additionally, each latent only attends to the bytes contained within its corresponding patch. In this setup, each patch $p_j$ is associated with a query vector. This vector is formed by aggregating the representations of the bytes within patch $p_j$, subsequently undergoing a linear transformation parameterized by $\mathcal{E}_{C} \in \mathbb{R}^{h_{\mathcal{E}} \times (h_{\mathcal{E}} \times U_{\mathcal{E}})}$, with $U_{\mathcal{E}}$ denoting the number of encoder cross-attention heads. Let $f_{\text{bytes}}(p_j)$ represent the sequence of bytes for patch $p_j$; under these definitions, we compute

\begin{align}
P_{0,j} &= \mathcal{E}_{C}(f_\text{bytes}((p_j)), f~ \text{is a pooling function} \\
P_l &= P_{l-1} + W_o\left(\text{softmax}\left(\frac{QK^T}{\sqrt{d_k}}\right)V\right) \\
\text{where } Q_j &= W_q(P_{l-1,j}), K_i = W_k(h_{l-1,i}), V_i = W_v(h_{l-1,i}) \\
h_l &= \text{Encoder-Transformer-Layer}_l(h_{l-1})
\end{align}

In this context, $P \in \mathbb{R}^{n_p \times h_{\mathcal{G}}}$ denotes the representations for $n_p$ patches, prepared for processing by the global model. The initialization of these representations is achieved by aggregating the byte embeddings $e_i$ associated with each patch $p_j$. The matrices $W_q$, $W_k$, $W_v$, and $W_o$ are utilized for computing the projections that define queries, keys, values, and outputs respectively. Here, keys and values are derived by projecting the byte representations $h_i$ produced by the prior layer (or the initial embeddings $e_i$ in the first layer). A patch-specific masking procedure is implemented so that each query $Q_j$ exclusively attends to keys and values associated with the bytes within its corresponding patch $j$. Multi-headed attention is applied to $Q$, $K$, and $V$, and since patch representations generally have higher dimensionality ($h_{\mathcal{G}}$) than $h_{\mathcal{E}}$, the patch representations $P_l$ are organized into multiple attention heads of size $h_{\mathcal{E}}$ during cross-attention, after which these heads are concatenated to reconstruct a vector of dimension $h_{\mathcal{G}}$. Furthermore, pre-LayerNorm is applied to queries, keys, and values, while positional embeddings are omitted from this cross-attention layer. A residual connection also envelops the cross-attention component.

\subsection{Local Decoder} 

The local decoder $\mathcal{D}$, much like the local encoder, utilizes a transformer architecture but with significantly fewer layers ($l_{\mathcal{D}} << l_{\mathcal{G}}$), making it lightweight. Its purpose is to convert the series of global patch representations $o_j$ back into the original raw bytes $y_i$. This process is autoregressive: the local decoder generates each raw byte conditioned on all previously outputted bytes. To do so, it leverages the hidden states outputted by the local encoder for the sequence of bytes. The architecture comprises $l_{\mathcal{D}}$ layers in which cross-attention and transformer layers alternate. Notably, each decoding block first applies a cross-attention operation to map from patch-level representations to byte-level ones before passing through the transformer's own layer, which subsequently refines these byte representations.

\subsubsection{Decoder Multi-headed Cross-Attention} 

In the decoder cross-attention, the roles of the queries and key/values are interchanged i.e. the byte-representations are now the queries, and the patch representations are now the key/values. The initial byte-representations for the cross-attention are initialized as the byte embeddings from the last encoder layer i.e. $h_{l_{\mathcal{E}}}$. The subsequent byte-representations for layer $l$, $d_{l,i}$ are computed as:

\begin{align}
D_0 &= h_{l_{\mathcal{E}}} \\
B_l &= D_{l-1} + W_o\left(\text{softmax}\left(\frac{QK^T}{\sqrt{d_k}}\right)V\right), \\
  &\text{where } Q_i = W_q(d_{l-1,i}), K_i = W_k(\mathcal{D}_C(o_j)), V_i = W_v(\mathcal{D}_C(o_j)) \\
  D_l &= \text{Decoder-Transformer-layer}_l(B_l)
\end{align}

In this formulation, the key and value projection matrices, $W_k$ and $W_v$, act upon the final patch representations $o_j$ from the global model, employing a linear transformation and the split operation $\mathcal{D}_C$. The query projection matrix, $W_q$, transforms the byte representations $d_{l-1}$ from the previous decoder transformer layer, or $h_{l_{\mathcal{E}}}$ if operating in the first layer. The result is then processed via the output projection matrix $W_o$, producing $B \in \mathbb{R}^{h_{\mathcal{D}} \times n_b}$, where $n_b$ designates the number of output bytes. Subsequently, the following decoder layer representations $D_l$ are produced by feeding the output from the cross-attention block, $B$, into a decoder transformer layer. Mirroring the approach of the local encoder cross-attention, this setup employs multi-head attention, applies pre LayerNorm, excludes positional embeddings, and integrates a residual connection encasing the cross-attention module.

\section{Experimental Setup}
\label{section:experiments}
We construct rigorous and controlled experimental conditions to assess the performance of {\textsc BLT} against models that utilize tokenization, ensuring that any potential benefits related to {\textsc BLT} accessing longer input contexts do not bias the comparison.

\subsection{Pre-training Datasets} For all model sizes investigated in this study, pre-training was conducted using two principal datasets: 1) The Llama 2 dataset~\citep{touvron2023llama}, containing 2 trillion tokens sourced from a diverse array of publicly accessible materials, subsequently subjected to rigorous cleaning and filtering protocols to ensure higher quality; and 2) {\textsc BLT}-1T, a new corpus comprising 1 trillion tokens extracted from assorted public resources, inclusive of a subset of the pre-training content provided by Datacomp-LM~\citep{li2024datacomp}. The Llama 2 dataset was utilized for scaling law analysis based on optimal token counts as identified by~\cite{dubey2024llama}, thereby guiding the selection of appropriate architecture for {\textsc BLT}. In contrast, {\textsc BLT}-1T was employed in a full-scale pre-training process to facilitate direct comparisons with Llama 3 across a range of downstream evaluation tasks. Importantly, neither dataset contains material collected from Meta’s products or services. For baseline trials involving tokenizer-based architectures, we adopted the Llama 3 tokenizer with a 128K vocabulary, as it yielded better performance than the Llama 2 tokenizer according to our empirical results.

\subsection{Entropy Model} For our experiments, the entropy model is implemented as a byte-level language model, trained on the identical data distribution as the comprehensive {\textsc BLT} framework. By default, this model is instantiated as a transformer architecture comprising 100M parameters, organized across 14 layers, with a hidden state size of 512 and utilizing sliding window attention over sequences of 512 bytes. All other hyperparameter values are aligned with those employed in both our local and global transformer models. We systematically evaluated variations in model scale, receptive field length, and architectural design, as detailed in \autoref{section:ablations}. Notably, when the receptive field is sufficiently narrow, the resulting entropy model can be represented compactly via a lookup table, yielding greater encoding efficiency.

\subsection{Entropy Threshold and Equalizing Context Length} For models employing entropy-based patching, we determine a patching threshold that yields a targeted mean \textit{patch size} across the pretraining data mixture. In {\textsc BLT}, the \textit{patch size} is not constrained by token boundaries and can be set to any value, which notably affects the amount of context provided to the model. To prevent models with larger patches from benefiting from increased context, we control for context length by keeping the expected number of bytes per batch fixed. Consequently, when larger patch sizes are used, we shorten the sequence length accordingly. Specifically, on Llama 2 data, a context window of 8k bytes is utilized, whereas on the {\textsc BLT}-1T dataset, the context window is raised to an average of 16k bytes, ensuring the batch size stays consistent at approximately 16M bytes.

Although the mean batch size remains unchanged, dynamic patching techniques result in varying proportions of bytes to patches during data batching. To optimize computational efficiency, our {\textsc BLT} training workflow aggregates batches based on patch count, thereby bypassing the need for additional padding within the costly latent transformer module. This approach guarantees that all batches comprise an identical quantity of patches. In the training phase, byte sequences are either padded or truncated—capped at 12k bytes for Llama 2 and at 24k bytes for the {\textsc BLT}-1T datasets—in order to mitigate memory usage spikes caused by the occasional presence of exceptionally large patches.

\subsection{Entropy Model Context}
\label{section:entropy-context}
Our observations indicate that applying entropy patching leads to increasingly larger patches in highly structured tasks such as multiple choice assessments, which are typified by frequent repetition (for example, refer to the patching illustration in an MMLU task in \autoref{fig:mmlu_patching}). The underlying factor behind this pattern is the reduced entropy associated with the duplicated content present in the entropy model's context window. Accordingly, in our extensive experiment with {\textsc BLT}-Entropy using a patch size of 4.5, we refresh the entropy context with line breaks and adopt an approximate monotonicity constraint; this approach is more resilient to "entropy drift" resulting from variable context lengths. It should be noted that this modification only changes the entropy calculation process; our method for determining the entropy threshold remains consistent.

\subsection{FLOPs Estimation} 
\label{section:flops}

Our methodology for calculating transformer FLOPs closely aligns with the approach described in Chinchilla~\citep{hoffmann2022training}, which encompasses the FLOPs associated with the feed-forward modules, qkvo projections within self-attention, as well as both the attention calculation and the output projection. One key deviation in our approach is the treatment of the input embedding layer: we consider it to be realized as an efficient lookup table rather than as a dense matrix multiplication, resulting in its contribution to the total FLOPs being negligible (0 FLOPs). In line with established practice, we further approximate that the backward propagation phase incurs a computational cost double that of the forward phase.

To compute flops \textit{per byte} for {\textsc BLT} models, we add up the flops for the local encoder transformer, the global latent transformer, and the local decoder transformer, together with the cross attention blocks in the encoder and the decoder:

\begin{align}
    \text{FL}_{\text{{\textsc BLT}}} &= \text{Transf. FL}(h_{\mathcal{G}}, l_{\mathcal{G}}, m=n_{ctx}/n_p,V=0) / n_p\\
   &+\text{Transf. FL}(h_{\mathcal{E}}, l_{\mathcal{E}}, m=w_{\mathcal{E}}, V=0) \\
   &+ \text{Transf. FL}(h_{\mathcal{D}}, l_{\mathcal{D}}, m=w_{\mathcal{D}}, V=256) \\ 
    &+ \text{Cross Attn. FL}(h_{\mathcal{E}}, l_{\mathcal{E}}, m=n_p, r=n_p/k) \cdot k/n_p \\ 
   &+ \text{Cross Attn. FL}(h_{\mathcal{D}}, l_{\mathcal{D}}, m=k, r=k/n_p) 
\end{align}

In this context, $n_{ctx}$ denotes the input sequence length measured in bytes, while $n_p$ represents the patch size parameter. The parameter $r$ stands for the proportion of queries to key/value pairs, and $k$ signifies the ratio between patch and byte dimensions, effectively describing how many local model partitions are concatenated to generate the overall model embedding (with $k=2$ depicted in \autoref{fig:crossattn}). The symbol $V$ designates the vocabulary size relevant for the output projection, which is exclusively employed within the local decoder component. Depending on the sequence type—whether it is based on bytes or patches—the attention mechanism utilizes a different context length denoted as $m$. We have adjusted the computation of attention-related flops individually for each architectural module. Detailed formulations for the FLOPs calculation of both Transformer and Cross-Attention operations can be found in Appendix~\ref{section:flops-appendix}.

\subsection{Bits-Per-Byte Estimation} Perplexity only makes sense in the context of a fixed tokenizer as it is a measure of the uncertainty for each token. When comparing byte and token-level models, following previous work~\citep{xue2022byt5, yu2023megabyte, wang2024mambabyte}, we instead report Bits-Per-Byte (BPB), a tokenizer independent version of perplexity. Specifically:

\begin{align}
    \text{BPB}(x)&= \frac{\mathcal{L}_{CE}(\pmb{x})}{\ln(2)\cdot n_{\text{bytes}}}
\end{align}

In this formulation, the total cross-entropy loss, which quantifies the uncertainty associated with the data $\pmb{x}$, is divided by both the total count of bytes present in $\pmb{x}$ and a normalization constant.

\subsection{Transformer Architecture Hyperparameters}

For every transformer block within {\textsc BLT}, encompassing both the local and global components, we generally adopt the Llama~3~\citep{dubey2024llama} architectural framework. In particular, feed-forward modules employ the SwiGLU activation function~\citep{swiglu}, and rotary positional embeddings (RoPE)~\citep{RoPe} with a rotation parameter of $\theta=500000$~\citep{xiong2024effective} are utilized exclusively within self-attention modules. We perform layer normalization via RMSNorm~\citep{RMSNorm}. For self-attention computations that leverage static attention patterns—such as \textit{block causal} or \textit{fixed-window block causal} masks—Flash attention~\citep{dao2022flashattention} is deployed, using a window size of 512 for fixed-width masks. As the cross-attention layers make use of dynamic masks based on patch-level dependencies, Flex Attention\footnote{\url{https://pytorch.org/blog/flexattention}} is leveraged to implement efficient fused kernels, thereby notably accelerating the training process.

\subsection{{\textsc BLT}-Specific Hyperparameters} 

To evaluate the performance of {\textsc BLT} models, we perform experiments in two key areas: assessing scaling patterns and examining downstream task performance. For these experiments, we include models with varying capacities: 400M, 1B, 2B, 4B, and 8B parameters. The architectural settings for these models are detailed in Appendix~Table~\ref{tab:arch}.

For the local encoder, max-pooling is used to generate the initial queries for the first cross-attention layer. Across all {\textsc BLT} variants, we set the number of hashes to $500,000$ and employ a single hash function, with n-grams spanning lengths from 3 to 8. Each model utilizes a uniform learning rate of $4e-4$. Our learning rate selection for both token and {\textsc BLT} models is grounded in a hyperparameter search within the interval $1e-3$ to $1e-4$, where experiments at the 400M and 1B model sizes indicated the identical learning rate as optimal.

For scaling law analyses using Llama-2 data, training batch-sizes are adopted as recommended by~\cite{dubey2024llama}, or, where appropriate, an equivalent byte-based measure. Model optimization leverages the AdamW optimizer~\citep{adamw} with $\beta_1 = 0.9$, $\beta_2 = 0.95$, and $\epsilon = 10^{-8}$. We apply a linear learning rate warm-up over 2000 steps, after which a cosine decay brings the learning rate down to zero. Weight decay is set at 0.1, and we employ global gradient clipping with a cutoff at 1.0.

\section{Scaling Trends}
\label{section:scaling}

We offer a comprehensive analysis of how byte-level models scale, providing insights that may guide future development of {\textsc BLT} models. Our investigation seeks to overcome the shortcomings of earlier studies on byte-level models in several key respects: (a) By examining performance patterns under the regime where compute usage is optimized for training, (b) By training equivalent 8B-parameter models using substantial training corpora—reaching up to 1T tokens or 4T bytes—and assessing their effectiveness on various downstream benchmarks, and (c) By exploring scaling dynamics when inference costs are kept consistent. In subsequent sections, we delve into particular benefits that stem from processing byte-based sequences.

\subsection{Parameter Matched Compute Optimal Scaling Trends}
\label{section:parameter-matched-scaling}
We conduct experiments on the Llama 2 dataset by training a series of \textit{compute-optimal} bpe and {\textsc BLT} models at four distinct scales, spanning from 1B to 8B parameters. Training flops are mapped against language modeling accuracy on a representative subset of the mixed training corpus. The bpe models adhere to the compute-optimal parameter-to-data ratios recommended by Llama~3~\citep{dubey2024llama}, following guidelines to maximize language modeling performance under a constrained computational budget~\citep{hoffmann2022training}. This approach furnishes a rigorous reference point for our analysis. Corresponding {\textsc BLT} models are trained in parallel for each bpe model, utilizing Latent Transformers with matched architecture and parameter counts to ensure fair comparison.

As shown in~\autoref{fig:scaling-compute-opt} (right), {\textsc BLT} models consistently equal or surpass the performance of their bpe equivalents, and this relationship persists across increasing model sizes and floating point operation counts. Based on existing literature, {\textsc BLT} represents the first byte-level Transformer architecture to exhibit scaling behavior on par with BPE-based models under compute-optimal conditions. This result supports our hypothesis that the ideal parameter-to-compute ratio for bpe also extends to {\textsc BLT}, or is at minimum closely aligned.

Both enhancements to model architecture and the adoption of dynamic patching are essential for achieving scaling performance comparable to bpe-based methods. As shown in ~\autoref{fig:scaling-compute-opt} (left), we present a comparison between space-patching-centric approaches and Llama 3. To estimate the performance of SpaceByte \citep{slagle2024spacebyte}, we utilize a version of {\textsc BLT} that applies space-patching, omitting n-gram embeddings and cross-attention mechanisms. While SpaceByte demonstrates improvements relative to Megabyte, there remains a notable gap compared to Llama 3. Turning to \autoref{fig:scaling-compute-opt} (right), we depict the performance gains attributable to both structural modifications and the integration of dynamic patching. At scale, {\textsc BLT} models achieve results competitive with leading tokenizer-based architectures like Llama 3.

Additionally, we note that for tokenizer-based models, selecting the Llama-3 tokenizer leads to superior performance compared to models utilizing the Llama-2 tokenizer, even when both are trained on identical datasets.

In conclusion, our {\textsc BLT} model exhibits scaling behaviors intermediate between Llama 2 and 3 when employing notably larger patch sizes. While the bpe tokenizers for Llama 2 and 3 demonstrate average token lengths of 3.7 and 4.4 bytes, respectively, {\textsc BLT} attains analogous scaling behavior when the average patch size increases to 6 or even 8 bytes. Because the inference flop count decreases as the average patch size grows, selecting an 8-byte patch size can deliver nearly 50\% reduction in inference computation. Additionally, models configured with larger patch sizes appear to improve their performance as the scale of model parameters and dataset size increases. Specifically, a {\textsc BLT} model with an 8-byte patch size initially underperforms bpe-tokenized Llama 2 at the 1B parameter mark, yet surpasses its performance at the 7B scale. These observations imply that larger patch sizes may yield further benefits as models and data sets grow, and that utilizing even more substantial patch sizes could become practical with future increases in model size and computational resources.

\subsection{Beyond Compute Optimal Task Evaluations}
\label{section:evals}
To further investigate scaling behavior, we train an 8B {\textsc BLT} model past the compute-optimal threshold using the expanded {\textsc BLT}-1T dataset, which offers a broader and higher-quality data corpus. The model’s abilities are then assessed across a diverse set of established benchmarks for both classification and generation tasks. The evaluation set comprises tasks aimed at measuring common sense reasoning, general knowledge, and code generation skills:
\paragraph{Classification tasks} We include ARC-Easy (0-shot)~\citep{Eval_ARC}, Arc-Challenge (0-shot)~\citep{Eval_ARC}, HellaSwag (0-shot)~\citep{Eval_hellaswag}, PIQA (0-shot)~\citep{Eval_piqa}, and MMLU (5-shot)~\citep{hendrycks2020measuring}. Performance is determined by prompt-scoring, where we compute the likelihood assigned to each possible choice and average the resulting accuracy rates.
\paragraph{Coding related generation tasks:} For measuring code synthesis in LLMs, we present pass@1 metrics for MBPP (3-shot)~\citep{Eval_MBPP} and HumanEval (0-shot)~\citep{Eval_HumanEval}, reflecting the model's proficiency in generating correct Python code snippets.

In \autoref{tab:evals}, we present a comparison among three models, each trained on the {\textsc BLT}-1T dataset: a model utilizing the bpe Llama 3 tokenizer,\footnote{The Llama 3 tokenizer, with its 128k vocabulary size, is selected because it outperforms the Llama 2 tokenizer, which has a 32k vocabulary.} as well as two different implementations of the {\textsc BLT} architecture. The first variant incorporates a space-patching method ({\textsc BLT}-Space), while the second relies on an entropy-driven patching strategy ({\textsc BLT}-Entropy), both applying an approximate monotonicity constraint and resetting the entropy-based model’s context at new line boundaries (further detailed in~\autoref{section:entropy-context}). All models are trained using an identical amount of computational resources (flop budget). Notably, for {\textsc BLT}-Entropy, we further refine performance by lowering the entropy threshold at inference from 0.6 to 0.1, which enhances task results albeit with an increased number of inference steps.

The {\textsc BLT}-Entropy model achieves superior results compared to the Llama 3 model on four of the seven evaluated tasks, even though both were trained on an equal volume of data. This performance gain can likely be attributed to two primary factors: (1) more efficient utilization of training computation made possible by dynamic patching, and (2) the approach of modeling information directly at the byte level rather than at the token level.

Conversely, {\textsc BLT}-Space lags behind the Llama 3 tokenizer in every task except one, yet its use of a larger average patch size of 6 bytes leads to a notable decrease in inference flops. By contrast, models based on bpe and entropy-patching display similar average patch sizes, each around 4.5 bytes on the mixed training dataset. Given equivalent training resources, the model employing larger patches is able to process 30\% more data than the bpe and entropy-patching counterparts. However, this increased data coverage may cause {\textsc BLT} to deviate further from achieving compute-optimality.

\subsection{Patches Scale Better Than Tokens} {\textsc BLT} models allow for joint scaling of both model capacity and patch size, all while keeping the computational budget for both training and inference fixed and the training dataset unchanged. The capability to enlarge patch size without bound is exclusive to patch-based architectures, liberating them from the efficiency limitations characteristic of models built on fixed-vocabulary tokens, as outlined in Section~\ref{section:bpe-patch}. Utilizing larger patches reduces the frequency of computation, which in turn enables the allocation of saved computational resources to increase the capacity of the global latent transformer, since it is invoked less frequently.

To evaluate whether increasing model size while reducing the number of steps on larger patches yields superior performance compared to training smaller models with more steps, we carry out a fixed inference scaling analysis. Utilizing models of 400m and 3.6B parameters equipped with the Llama 2 tokenizer as a baseline, we identify flop-matched models using the Llama 3 tokenizer, in addition to {\textsc BLT}-Entropy models employing mean patch sizes of 6 and 8 bytes on the training data mixture (refer to~\autoref{tab:fixed-inf-params} for specifics about model configurations). For the models with a patch size of 8, we increase the number of encoder layers from 1 to 3. Across all configurations, we train the models under different training flop budget constraints.

As shown in \autoref{fig:fixed_inference_scaling}, {\textsc BLT} models demonstrate superior scaling behavior compared to tokenization-based models across both categories of inference flops. For small training resources, BPE models initially exhibit stronger performance; however, once beyond the compute-optimal point, {\textsc BLT} quickly outperforms them. In many practical scenarios, it may be advantageous to allocate greater resources to pretraining in a one-time investment, resulting in a model that achieves higher performance under a constrained inference budget. The 8B parameter class of models, such as Llama 3.1, exemplifies this strategy, having been exposed to a training corpus orders of magnitude larger than what would be considered optimal based solely on compute for that scale.

The point at which {\textsc BLT} surpasses token-based models now occurs at a lower multiple of the compute-optimal budget when advancing to models in the higher flop class, shifting from three times down to 2.5 times the optimal budget. Additionally, among the higher flop class models, increasing the patch size to 8 results in a sharper scaling curve, enabling it to outperform the other models at a smaller scale. As outlined in Section~\ref{section:parameter-matched-scaling}, using larger patch sizes brings performance closer to that of BPE models as model size grows. We partly attribute this trend to the diminishing proportion of total computation allocated to the byte-level Encoder and Decoder, which seem to exhibit slower scaling than the Latent Transformer. Notably, when the total parameter count is scaled by a factor of 20, from 400M up to 8B, the increase in {\textsc BLT}’s local model parameters is only about twofold. This distinction matters because higher patch sizes only influence flops consumed by the patch Latent Transformer, without affecting the byte-level modules. As a result, the {\textsc BLT}-Entropy model with patch size 8 increased from 1.6 times to 1.7 times the size of the Llama 2 model in transitioning to the higher model scale.

To summarize, our investigation into patch-length scaling reveals that the {\textsc BLT} patch-based framework exhibits more favorable scaling properties when both patch size and overall model size are increased in tandem. These advantageous scaling behaviors are sustained and even enhanced as the model size grows further.

\section{Byte Modeling Improves Robustness} We further evaluate how {\textsc BLT}'s robustness compares to that of token-based models which do not access byte-level representations, and introduce a method to incorporate byte-level processing into pretrained token-based architectures.
\label{sec:robustness}

\subsection{Character-Level Tasks}

One of the initial incentives for developing byte-level models stemmed from their ability to handle byte-level noise within inputs, as well as their capacity to recognize the substructures that compose tokens—a challenge for present-day models reliant on tokenization. To assess these aspects, we conduct supplementary experiments on tasks that probe the models’ robustness against noisy inputs and their sensitivity to character-level information across both English and multiple languages, encompassing numerical digits and phonemes. The outcomes of these evaluations are provided in Table~\ref{tab:char_tasks}.

\paragraph{Noisy Data} To assess the robustness of tokenizer-based models in comparison to {\textsc BLT}, we generate corrupted variants of the benchmark classification tasks outlined in Section~\ref{section:evals}. Specifically, we implement five different character-level noising techniques to introduce distortions into the text: (a) \textit{AntSpeak}: Transforms the text entirely into uppercase letters, inserting spaces between each character. (b) \textit{Drop}: Randomly deletes 10\% of the characters in the input text. (c) \textit{RandomCase}: Randomly assigns half of all characters to uppercase and the remaining half to lowercase. (d) \textit{Repeat}: Selects 20\% of the characters at random and repeats each one as many as four times. (e) \textit{UpperCase}: Converts all text characters to uppercase. In our evaluation protocol, each noising strategy is independently applied to the prompt, to the completion, and to both, yielding distinct evaluation scenarios. We then present mean performance scores for these settings. As detailed in Table \ref{tab:char_tasks}, results on noisy versions of the HellaSwag dataset~\citep{Eval_hellaswag} demonstrate that {\textsc BLT} consistently surpasses tokenizer-based models in resilience to textual distortions, offering, on average, an 8-point improvement over a similarly trained baseline model, and even outperforming the Llama 3.1 model, which is trained on a significantly larger corpus.

\paragraph{Phonology - Grapheme-to-Phoneme (G2P)} We evaluate {\textsc BLT}'s proficiency in converting a string of graphemes—letters that constitute a word—into its corresponding phonemic representation, capturing its pronunciation. Table \ref{tab:char_tasks} displays the outcomes of the G2P experiment under a 5-shot regime, utilizing the Phonology Bench~\citep{suvarna2024phonologybench} dataset. The findings indicate that {\textsc BLT} achieves superior performance relative to the baseline model equipped with the Llama 3 1T tokenizer.

\paragraph{CUTE}
To evaluate character-level comprehension, we measure the performance of {\textsc BLT} on the CUTE benchmark~\citep{edman2024cute}, which contains a range of tasks grouped into three main types: compositional understanding, orthographic similarity analysis, and sequence manipulation. This suite of tasks is notably demanding for most models utilizing tokenizers, as such models typically recognize their token spellings but have difficulty leveraging this knowledge for text manipulation. According to Table~\ref{tab:char_tasks}, {\textsc BLT}-Entropy surpasses both BPE Llama 3 variants on CUTE by a margin exceeding 25 points. Notably, our model excels in character manipulation tasks, attaining 99.9\% accuracy on both spelling-related subtasks. These substantial gains, achieved even though {\textsc BLT} was exposed to 16 times less training data than Llama 3.1, suggest that BPE-based approaches face inherent challenges with acquiring character-level skills. Figure~\ref{fig:cute_examples} depicts cases in which the Llama 3 tokenizer-based system fails while the {\textsc BLT} model is successful. The BPE approach does outperform on word deletion and insertion; these forms of word manipulation are less accessible to a byte-level model, though the difference remains modest. Moreover, the transition from recognizing characters to understanding words may ultimately prove less complex than the reverse. For all assessments, we adhere to the identical evaluation protocol and employ the original Huggingface prompts. Additional prompt design could further enhance the outcomes for BPE-based models.

\paragraph{Low Resource Machine Translation} We assess {\textsc BLT} on translation tasks involving six major language families and twenty-one lower-resource languages, covering a diverse range of scripts, using the FLORES-101 benchmark~\citep{goyal2022flores}. SentencePiece BLEU scores are presented in Table~\ref{tab:flores}. The findings indicate that {\textsc BLT} surpasses a model utilizing the Llama 3 tokenizer, showing an overall improvement of 2 BLEU points when translating into English and a 0.5-point gain when translating from English. In widely used language pairs, {\textsc BLT} matches or slightly exceeds the performance of Llama 3. Notably, for many language pairs in lower-resource families, {\textsc BLT} substantially outperforms Llama 3, highlighting the advantages of byte-level modeling in capturing and generalizing over rare and unique byte sequences.

\subsection{Training {\textsc BLT} Utilizing Llama 3}
\label{sec:distillation}
We investigate an approach in which {\textsc BLT} models benefit from pre-existing, pretrained tokenizer-based models to accelerate and stabilize the training process. In this method, the global transformer parameters in {\textsc BLT} are initialized with the weights from a Llama 3.1 checkpoint. The subsequent training procedure involves updating these global parameters using a learning rate set to one-tenth of that used for the local encoder and decoder modules, while keeping the total number of optimization steps consistent with Llama 3's recommended setting. Comparative results with a baseline {\textsc BLT} model are summarized in Table~\ref{tab:distill}. The data indicates that the {\textsc BLT} variant initialized from Llama 3.1 exhibits superior performance relative to both the original Llama 3 and the standard {\textsc BLT} baseline, despite all being trained with equivalent computational budgets. Furthermore, in relation to our {\textsc BLT}-Entropy variant (see Table~\ref{tab:evals}), which was exposed to a much larger corpus of 1T tokens, {\textsc BLT} initialized from Llama 3.1 still surpasses it on the MMLU benchmark, highlighting the efficiency of this approach for achieving competitive results while substantially reducing the required training flops.

This approach can alternatively be understood as reconfiguring models reliant on tokenizers into ones that do not require them, thereby adapting the pre-trained LLaMA 3.1 model into a {\textsc BLT} model. For a thorough evaluation, Table \ref{tab:distill} presents results from the original LLaMA 3.1 model, which was trained on 15T tokens, and compares its performance to that of the {\textsc BLT} model based on LLaMA 3. We observe that while the transition results in marginal decreases in MMLU and HumanEval scores, the performance on several other benchmarks deteriorates more noticeably. This indicates the need for continued research to maximize the capabilities of the pre-trained model, particularly by refining data composition strategies and adjusting other crucial hyperparameters.

\section{Ablations and Discussion}
\label{section:ablations}
Here, we analyze the results of ablation studies to support the architectural decisions made for {\textsc BLT}. We further examine the patching approach and the selection of hyper-parameters implemented in the 8B parameter {\textsc BLT} model trained using the {\textsc BLT}-1T dataset.

\paragraph{Entropy Model Hyper-parameters} To investigate how the scale of entropy models and the length of context windows impact model scaling, we develop a range of byte-level entropy transformer models with sizes spanning from 1m to 100m parameters, paired with context window lengths varying from 64 up to 512. To assess scaling laws, we plot bpb against training FLOPs using our $400m$ and $1b$ {\textsc BLT} models, both of which were trained on the Llama-2 dataset; results are displayed in \autoref{fig:entropy_ablation}. Our findings demonstrate that improvements in scaling performance are associated with increases in both model size and context window length, although we observe diminishing gains when surpassing 50m parameters.

\paragraph{Types of Patching}

We conduct an ablation study on the four patching schemes outlined in Section~\ref{section:patching}, namely: 1) Strided Patching employing strides of 4 and 6, 2) Patching at whitespace positions, 3) BPE Tokenizer patching utilizing the Llama 3 tokenizer, and 4) Entropy-based patching facilitated by a compact byte LLM.

Although dynamic patching shortens the actual sequence lengths, we regulate the sequence length to ensure that every patching strategy preserves an equivalent context window. Across both training and inference, each model is exposed to the same expected number of bytes per sequence, thereby eliminating potential confounds related to the capacity for modeling longer contexts. The results of these controlled ablation studies are presented in Figure~\ref{fig:scaling-compute-opt}. All alternative patching approaches show better performance than static patching, with space patching yielding results nearly indistinguishable from those of dynamic entropy based patching.

\autoref{tab:evals_ablation} displays benchmark results for {\textsc BLT} models, examining the performance of tokenizer-based approaches, space patching, and entropy-based patching. All models were trained on the Llama 2 dataset for an empirically determined number of training steps~\citep{dubey2024llama}. While space patching offers a more straightforward method that avoids online entropy estimation during training, our experiments confirm that the improvements attributed to entropy-based patching observed in scaling studies~(Section~\ref{section:scaling}) are retained in downstream evaluation benchmarks as well.\footnote{Space patching results correspond to initial experiments conducted without cross-attention; nevertheless, similar patterns persist when cross-attention is included.}

\paragraph{Cross-Attention}
\autoref{tab:crossattn_ablations_1b} presents results from ablation studies evaluating the placement of cross-attention mechanisms at distinct stages within the encoder and decoder of {\textsc BLT}. Regarding cross-attention in the encoder, we compare three methods for initializing the queries: (1) assigning all global states a uniform learned embedding, (2) employing a hash-based embedding derived from the patch’s bytes, and (3) utilizing a pooled version of the encoder’s hidden states corresponding to the patch bytes at the specific encoder layer.

Our results indicate that applying cross-attention within the \textit{decoder} yields the greatest benefits. While integrating cross-attention in the encoder leads to marginal gains, these improvements occur exclusively when queries are initialized via pooling. Furthermore, cross-attention proves particularly advantageous for Common-Crawl data and is especially effective when utilizing larger patch sizes.

\paragraph{n-gram Hash Embeddings} 

We experiment with n-gram hash embedding vocabulary sizes of 0, 100K, 200K, and 400K, with the outcomes reported in Table~\ref{tab:ngram_ablations_1b}. Our observations suggest that hash embeddings offer improvements across all evaluated domains, with the most pronounced effects seen on Wikipedia and Github (showing a 0.04 bpb difference, compared to a 0.01 bpb difference after 15k steps at the 8B scale). At the 8B parameter setting, reducing the vocabulary size from 500K to 300K hashes results in only a minor change—approximately 0.001 bpb after 15k steps. These findings imply that incorporating hashes is crucial for achieving {\textsc BLT} performance on par with tokenizer-based approaches. Nonetheless, beyond 300K hashes, the benefit of adding more hashes diminishes. Furthermore, the enhancements from hash embeddings seem to act in a complementary fashion to cross-attention mechanisms, as the performance gains occur on distinct datasets.

\paragraph{Local Model Hyperparameters} Table \ref{tab:local_layers} presents an ablation study on different configurations for the layer counts in both the local encoder and decoder. The results indicate that, when used with hash n-gram embeddings, {\textsc BLT} performs effectively even when the encoder is minimal—consisting of only a single layer—while the decoder benefits from a greater depth.

\section{Related Work}
\label{section:related_work}

\textbf{Character-Level RNNs:} The task of Character Language Modeling has remained prominent since the advent of neural models~\citep{sutskever2011generating,mikolov2012subword,graves2013generating}, primarily because these models naturally handle unseen words without reliance on back-off strategies. \cite{kim2016character} present an approach where input sequences are processed at the character level by convolutional and highway networks, which then interface with LSTM-based RNNs; their method achieves competitive results with contemporary RNN language models on English, and surpasses them when applied to languages with rich morphological structures—a notable benefit of character-level LLMs. In related work, \cite{kenter2018byte} demonstrate that byte-level LSTM architectures excel at machine comprehension tasks in morphologically complex languages such as Turkish and Russian, outperforming analogous word-based systems. Building on these ideas, \cite{zhang2015character} employ convolutional models at the character level for text classification and report superior accuracy in certain scenarios compared to word-based counterparts. To further advance performance in character-level language models, \citet{chung2019hierarchical} propose hierarchical LSTM architectures equipped with boundary detectors at different granularity levels to reveal underlying text structure. Additionally, ByteNet \cite{kalchbrenner2016neural} introduces the use of convolutional networks operating directly over characters, instead of attention mechanisms, for machine translation tasks.

\textbf{Character-Level Transformers:} The integration of attention mechanisms in transformer architectures~\citep{vaswani2017attention}, alongside the use of subword tokenization~\citep{sennrich-etal-2016-neural}, has led to notable gains in the performance of neural networks for both language modeling and standard evaluation tasks. Nevertheless, relying on word or subword units introduces a built-in inductive bias regarding the granularity at which these models operate. In an effort to leverage the strengths of transformers and capitalize on earlier encouraging outcomes for character-level language modeling, \citet{al2019character} employed extremely deep transformer architectures. By introducing auxiliary losses, they successfully trained transformer-based models that surpassed the performance of preceding LSTM-based character language models. Despite these improvements, their models still lagged noticeably behind those trained on word-level representations. Furthermore, GPT-2~\citep{radfordgpt2} reported that when working with expansive datasets such as the 1 Billion Word Benchmark, byte-level language models failed to achieve competitiveness with their word-level counterparts.

Although \cite{Choe2019BridgingTG} showed that transformer-based byte-level LLMs can surpass subword-level LLMs of similar size, these models demand significantly greater computational resources and lengthier training times. In a related line of work, \cite{el2020characterbert} introduced CharFormer, a variant of BERT employing convolutions over character embeddings to construct word-level representations; their approach yielded performance gains within the medical domain but at the expense of elevated compute requirements. Further, \cite{clark2022canine} presented CANINE—a 150M parameter, encoder-only model that processes raw character sequences. CANINE’s architecture consists of a deep transformer stack, paralleling the global model described in this work, combined with local transformers and strided convolutions that enable effective downsampling of input characters; the resulting model demonstrated superior downstream multilingual performance relative to a token-based encoder-only counterpart (mBERT). ByT5~\citep{xue2022byt5} examined strategies for byte-level encoder-decoder models eschewing any patching operations. These models displayed enhanced robustness to noisy inputs and remained competitive with traditional tokenizer-based models even with a data reduction by a factor of four; nevertheless, the omission of patching required computation of expensive attention mechanisms across all bytes, resulting in substantial computational overhead. Modeling inputs directly at the byte level, as opposed to subwords, invariably extends sequence length and poses scalability challenges. In recent developments, the Mamba architecture~\citep{gu2023mamba}, with its capacity to retain a fixed-size memory across extensive contexts, enabled \cite{wang2024mambabyte} to train a byte-level Mamba model—again, without patching—that surpassed byte-level transformer baselines, in flop-constrained scenarios at a 350M parameter scale, based on bits-per-byte across multiple datasets.

\textbf{Patching-based approaches:} Leveraging patching techniques can significantly reduce the otherwise high computational cost associated with byte-level LLMs, while potentially maintaining their performance. A variety of studies have reported promising outcomes in limited settings with respect to both model size and training data volume. For instance, \cite{nawrot-etal-2022-hierarchical} explored static patching strategies that incorporate downsampling and upsampling procedures, resulting in the introduction of the hourglass transformer, which surpassed alternative byte-level models at the 150M parameter scale. Building upon these ideas, \cite{nawrot-etal-2023-efficient} introduced enhancements through dynamic patching methods, incorporating boundary predictors trained in an end-to-end manner, as well as predictors trained with supervision from tokenizers, and an entropy-informed patching strategy similar to {\textsc BLT}. Their results demonstrated that such dynamic approaches could outperform standard transformers of the period on language modeling benchmarks, specifically with models containing approximately 40M parameters trained on around 400M tokens. In another line of research, \cite{lester2024training} considered training on data sequences compressed using arithmetic coding, enabling greater compression compared to BPE. Utilizing an equal-info windows technique, they exceeded the performance of byte-level models on language modeling tasks; however, their approach remained inferior to subword-level baselines.

Our research is primarily motivated by and bears close resemblance to MegaByte~\citep{yu2023megabyte}, a decoder-only causal LLM that implements a fixed, unchanging patching process and concatenates encoded representations to translate bytes into patches. On the decoder side, it employs a localized model to revert from patches back into bytes. The authors demonstrate that, at the 1B parameter scale and across approximately 400B bytes of data, MegaByte achieves comparable performance to models utilizing tokenizers. Throughout our study, we systematically compare MegaByte in all experimental settings and observe that, under identical FLOP budgets, static patching underperforms relative to contemporary compute-optimal tokenizer-based models. We further show how {\textsc BLT} effectively narrows this performance gap. In line with this, \cite{slagle2024spacebyte} document a similar finding regarding MegaByte. They propose enhancements such as expanding the static patching approach to encompass patching at whitespaces and similar space-like byte positions, and introduce a local encoder component. These modifications yield gains over tokenizer-based transformer architectures within certain compute constraints, notably on domains such as Github and arXiv when trained at the 1B parameter scale. Our evaluation of this model corroborates these outcomes and further indicates that substantial architectural innovations are required to allow byte-level models to scale and truly achieve parity with the latest token-based models like Llama 3.

\section{Limitations and Future Work}

For this study, when making architectural decisions, we train our models for the step counts that have been identified as optimal for Llama~3~\citep{dubey2024llama}. Nevertheless, these scaling laws were originally developed for BPE-level transformers and might produce non-optimal ratios of data to model parameters when applied to {\textsc BLT}. The precise determination of scaling laws specific to {\textsc BLT} is left as an open problem for future investigation, with the possibility that such research could yield more advantageous scaling dynamics for our architecture. Furthermore, the majority of our experiments have been restricted to model sizes up to 1B parameters; thus, it remains possible that as we move to larger models—on the order of 8B parameters or more—different architectural configurations could prove superior, potentially resulting in further gains at these greater scales.

Most current transformer codebases are optimized primarily for tokenizer-based architectures. Although our work includes theoretical flop-equivalent comparisons and leverages specialized efficient modules like FlexAttention to support components that differ from the standard transformer, our current implementations might still lag behind tokenizer-based models with respect to wall-clock efficiency. Further engineering and optimization could help close this performance gap.

Although {\textsc BLT} relies on a separately trained entropy model for the patching process, integrating the training of the patching model within an end-to-end framework presents a promising avenue for future research. Section \ref{sec:distillation} details preliminary experiments that demonstrate encouraging results for converting tokenizer-based models—such as Llama 3, which has been trained on over 10T tokens—into byte-based models. This is achieved by initializing and freezing the global transformer with the pretrained weights. Advancing this line of investigation may yield techniques that not only preserve the advantages offered by bytefying, but also enhance performance beyond that of existing tokenizer-based models, potentially eliminating the need for training entirely from the beginning.

\section{Conclusion}
\label{section:conclusion}

In this work, we introduce the Byte Latent Transformer (\textbf{BLT}), an innovative architecture that moves away from the traditional reliance on fixed-vocabulary tokenization in large language models. Through the implementation of a flexible, trainable approach to aggregating bytes into patches, {\textsc BLT} dynamically allocates computational effort according to the complexity of input data, yielding notable enhancements in both computational efficiency and model robustness. Our comprehensive scaling analysis reveals that {\textsc BLT} achieves parity with tokenization-dependent architectures such as Llama 3 for models up to 8B parameters trained on 4T bytes. Moreover, it is capable of exchanging modest declines in evaluation performance for inference flop reductions of as much as 50\%. In addition, {\textsc BLT} enables a novel pathway for model scaling by supporting concurrent enlargement of both model and patch sizes while maintaining a constant inference budget. This paradigm proves particularly effective within the computational limits typical of real-world deployments. When working directly on raw byte sequences, {\textsc BLT} further strengthens the model’s handling of infrequent data cases, resulting in greater resilience to noisy examples and more nuanced comprehension of sub-word patterns. Taken together, these findings establish {\textsc BLT} as a compelling substitute for existing tokenization-based models, offering a robust, flexible, and highly scalable solution for building next-generation language models.

\end{document}